\section{Introduction}
\label{sec:introduction}

Take a Bangla offensive-language classifier that ranks in-domain comments almost perfectly (AUROC 0.97). Show it the
same kind of comments typed in Latin script, as many Bangla speakers write online, and its false-negative control,
calibrated on its own data, misses more than four out of five toxic comments. Yet it still ranks those comments far
better than chance (AUROC about 0.8). The model has not forgotten what offensive language looks like; it has kept
the wrong cut-off. Temperature scaling \citep{guo2017calibration}, the default repair for miscalibrated confidence,
cannot help here: it rescales confidence without changing a single decision.

This paper asks what actually breaks when text classifiers meet the language varieties they were not trained on, and
what it takes to repair it. We fine-tune seven encoder families (five seeds each) on Bangla offensive-language and
sentiment data and compare them with four instruction-tuned LLMs on identical items. The test conditions cover regional
dialects, including a Sylheti set we translated, natural Bangla--English code-mixing written in two scripts,
tri-lingual Bangla--English--Hindi code-mixing, and other languages. To check that the conclusions are not an artefact
of one source language, we repeat the analysis with encoders trained on English (OLID) and on Hindi--English code-mixed
data (SentiMix Hinglish).

\paragraph{What we find.}
\begin{itemize}
  \item \textbf{Shift has a small number of regimes, and the common one is a threshold shift}
  (Section~\ref{sec:results-regimes}). Dialects need no correction. Temperature shifts are rare and small. For every
  source language, the dominant failure keeps the ranking and moves the threshold. A genuine collapse of the ranking
  occurs once.
  \item \textbf{Temperature is the wrong knob} (Section~\ref{sec:results-threshold}). Under a threshold shift, fitting
  the temperature on target labels lowers the log-loss but changes no decision, while a label-free bias correction
  raises macro-F1 by about 0.09. That correction silently assumes balanced classes and does harm when the target is
  skewed, as moderation targets are; a class prior estimated from 50 labels repairs it.
  \item \textbf{The collapse is about definitions, not language} (Section~\ref{sec:results-causes}). Changing only the
  script of identical comments shifts encoder AUROC by up to 0.22, but retraining the same architectures on the
  target's definition of offensiveness restores the collapsed ranking (pre-registered difference-in-differences
  $+0.378$, $p_{\mathrm{Holm}} < 0.001$).
  \item \textbf{Fifty labels go a long way} (Sections~\ref{sec:results-systems} and~\ref{sec:shift-triage}). Shift is
  easy to detect without labels but impossible to diagnose from the model's own confidence. About 50 target labels are
  enough to test whether a system's ranking survived, to choose between the fine-tuned encoder and an LLM, to correct
  the threshold and to certify coverage and false-negative rates after that choice.
\end{itemize}

Code, configurations, every derived table and a script that rebuilds all figures and tables of this paper are
released (Appendix~\ref{sec:appendix-reproducibility}).
